\chapter{Список источников и внешних ресурсов}
\label{ch:sources}
\index{источники}

\section{Нормативные и методические документы}

\begin{enumerate}
  \item ГОСТ~Р~72118-2025 --- требования к процессам жизненного цикла автоматизированных систем с учётом требований информационной безопасности (цели и предположения безопасности, доверенная вычислительная база, учёт компонентов). Официальный текст следует использовать в актуальной редакции.
  \item ISO/SAE~21434:2021 --- дорожная транспортная кибербезопасность (в т.\,ч. моделирование угроз и TARA). Официальное приобретение и применение --- по правилам владельца стандарта.
\end{enumerate}

\section{Учебные и справочные ресурсы в сети}

\begin{itemize}
  \item Пример изоляции компонентов и контролируемого взаимодействия (Jupyter-ноутбук, локальный учебный пример):\\
  \path{references/traffic_light_demo/cyberimmunity-traffic-lights-example.ipynb}
  \item Краткое руководство по примеру «Светофор»: \path{docs/traffic_light_demo.md}
\end{itemize}

Полный перечень команд и путей к файлам прототипа приведён в исходном репозитории учебного стенда; настоящий документ обобщает назначение систем без дублирования всех имён файлов.
